\documentclass[12pt]{article}
\usepackage[T1]{fontenc}
\usepackage[utf8]{inputenc}
\usepackage{lmodern}
\usepackage{textcomp}
\usepackage{enumitem}
\usepackage{geometry}
\geometry{margin=1in}
\begin{document}
\begin{center}
Answer Key - Philosophy - Graduate Level Assessment 20 \
\small (Answers are ordered exactly as in the question paper.)
\end{center}

\section*{Multiple Choice}
\begin{enumerate}[label=\arabic*.]
\item \textbf{Answer:} A
\\ \textit{Options:} A) logical justification.; B) epistemic justification.; C) existential justification.; D) theoretical justification.; E) moral justification.
\\ \textit{Note:} Gauthier argues that practical rationality is fundamentally rooted in logical justification, as it requires that actions and decisions be supported by coherent reasoning and principles, thereby linking rationality directly to the validity of the justifications that underpin them.
\item \textbf{Answer:} A
\\ \textit{Options:} A) are only applicable to certain cultures.; B) are determined by societal norms.; C) are only valid if they align with personal beliefs.; D) have no moral significance.; E) can be altered over time.
\\ \textit{Note:} Ross's theory of prima facie duties suggests that these moral obligations can vary in their application and significance depending on cultural contexts, highlighting the relativistic aspect of ethical principles.
\item \textbf{Answer:} D
\\ \textit{Options:} A) is merely a cultural artifact.; B) is not really relevant to our lives.; C) provides a comprehensive guide to conduct.; D) none of the above.
\\ \textit{Note:} Wolf argues that morality cannot be fully captured by any singular definition or theory, suggesting that the complexities of moral experience exceed the options provided in the question.
\item \textbf{Answer:} A
\\ \textit{Options:} A) proportional; B) equality; C) retaliatory; D) punitive; E) compensatory
\\ \textit{Note:} Nathanson supports proportional retributivism because he argues that the punishment should be proportionate to the severity of the crime committed, aligning with his philosophical stance that justice requires a balanced response to wrongdoing.
\item \textbf{Answer:} C
\\ \textit{Options:} A) Laudatory personality; B) Red herring; C) Reprehensible personality; D) Circular reasoning; E) Straw man fallacy
\\ \textit{Note:} The fallacy of a reprehensible personality occurs when an individual's negative traits are used to dismiss or negate their positive actions, thereby undermining the validity of their achievements based solely on their character flaws.
\end{enumerate}

\section*{True / False}
\begin{enumerate}[label=\arabic*.]
\setcounter{enumi}{5}
\item \textbf{Answer:} False
\\ \textit{Options:} A) True; B) False
\\ \textit{Note:} The natural law fallacy is not primarily a type of false cause because it involves the mistaken belief that what is natural is inherently good or just, rather than a misattribution of causality between unrelated phenomena.
\item \textbf{Answer:} True
\\ \textit{Options:} A) True; B) False
\\ \textit{Note:} The statement is true because affirming the consequent is a logical fallacy that occurs when one assumes that a specific outcome (membership) must be true based solely on the presence of evidence (denial of membership) that does not logically support that conclusion.
\item \textbf{Answer:} True
\\ \textit{Options:} A) True; B) False
\\ \textit{Note:} Singer's argument hinges on the ethical principle that complicity in wrongdoing, such as supporting corrupt regimes through business dealings, is morally equivalent to benefiting from stolen property, as both actions contribute to the perpetuation of injustice and exploitation.
\item \textbf{Answer:} True
\\ \textit{Options:} A) True; B) False
\\ \textit{Note:} Butler's philosophy emphasizes that justice is inherent in the moral framework that governs human interactions, positioning it as the foundational right that transcends other claims to ensure fairness and equality among individuals.
\item \textbf{Answer:} True
\\ \textit{Options:} A) True; B) False
\\ \textit{Note:} The statement is true because the fallacy of appeal to popularity, also known as argumentum ad populum, asserts that a belief or action is valid simply because it is widely accepted or endorsed by the majority, thereby undermining the logical basis for its justification.
\end{enumerate}

\section*{Long Form Response}
\begin{enumerate}[label=\arabic*.]
\setcounter{enumi}{10}
\item \textbf{Answer:} Attributing complex events to overly simplistic causes, often referred to as the fallacy of reductionism, undermines the nuanced understanding of causation in philosophical reasoning. This misattribution can lead to an incomplete analysis of events, obscuring the interplay of multiple contributing factors and the intricate web of relationships inherent in complex systems. By failing to recognize the multifaceted nature of causation, we risk oversimplifying moral responsibility and neglecting the contextual variables that shape outcomes. Such reductionist thinking not only distorts our comprehension of causative links but also diminishes our ability to engage meaningfully with ethical dilemmas and social phenomena, ultimately limiting philosophical discourse and understanding.
\\ \textit{Note:} correct because it emphasizes that reductionism simplifies complex events, leading to an inadequate analysis of causation that overlooks the interconnected factors involved, thereby distorting our understanding of moral responsibility and the complexity of systems.
\item \textbf{Answer:} De Marneffe's perspective on drug regulation presents a nuanced approach where drug use is deemed illegal, yet the sale and manufacturing of drugs are considered permissible. This stance implicates a moral framework that distinguishes between personal autonomy and social responsibility; individuals should not be allowed to engage in drug use due to the potential harm it poses to themselves and society, but the market dynamics of drug production and sale can be regulated to minimize illicit activity and promote safety. By advocating for the legality of sale and manufacturing, de Marneffe emphasizes the need for a controlled environment that can reduce the risks associated with drug use while acknowledging the economic realities of drug markets. Thus, this duality reflects a pragmatic approach to public health and safety, balancing individual rights with societal concerns.
\\ \textit{Note:} correct because it captures de Marneffe's ethical distinction between the protection of individual wellbeing through the illegality of drug use and the pragmatic regulation of the drug market to address societal harms while still allowing for legal sale and manufacturing.
\end{enumerate}

\vfill
\noindent\textit{End of Answer Key}
\end{document}